\chapter{Object-Oriented Programming (OOP)}

\textbf{Object-Oriented Programming (OOP)} is a programming paradigm based on the concept of "objects", which contain data (attributes) and code (methods). OOP enables modularity, encapsulation, and code reuse when building large software systems.

\section{Classes, Objects, \texttt{\textbackslash\{\}\_\textbackslash\{\}\_init\textbackslash\{\}\_\textbackslash\{\}\_}, and \texttt{self}}

A \textbf{Class} is a blueprint for creating objects. An \textbf{Object} is an instance of a class.

\begin{definitionbox}{Class Definition and Instance Attributes}
- \texttt{\_\_init\_\_(self, ...)}: The constructor method automatically executed when a new object instance is created.
- \texttt{self}: Represents the specific instance of the class being created or modified.
\begin{lstlisting}
class Student:
    def __init__(self, name, roll_no, gpa):
        self.name = name          # Instance attribute
        self.roll_no = roll_no
        self.gpa = gpa

    def display_info(self):       # Instance method
        return f"Student: {self.name} | GPA: {self.gpa}"

# Instantiating objects
s1 = Student("Alice", 101, 3.9)
print(s1.display_info())  # Output: Student: Alice | GPA: 3.9
\end{lstlisting}
\end{definitionbox}

\begin{videocard}{IITM BS Lecture: Introduction to OOP Basics}{https://www.youtube.com/watch?v=X1mOrCwZUgQ}
Official lecture introducing class concepts, objects, attributes, and methods.
\end{videocard}

\begin{videocard}{IITM BS Lecture: Attributes, Methods, Init, and Self}{https://www.youtube.com/watch?v=n_jPQlWhGhM}
Official lecture explaining the constructor \texttt{\_\_init\_\_} and instance binding via \texttt{self}.
\end{videocard}

\section{Inheritance and Method Overriding}

\textbf{Inheritance} allows a child class to inherit attributes and methods from a parent class, promoting code reuse.

\begin{theorembox}{Parent vs Child Classes}
A child class can override methods from the parent class to specialize behavior, and call \texttt{super().\_\_init\_\_()} to reuse parent initialization logic.
\end{theorembox}

\begin{examplebox}{Model Inheritance Architecture}
\begin{lstlisting}
class BaseModel:
    def __init__(self, model_name):
        self.model_name = model_name

    def fit(self, X, y):
        raise NotImplementedError("Subclasses must implement fit()")

class LinearRegressionModel(BaseModel):
    def __init__(self, learning_rate=0.01):
        super().__init__("Linear Regression")
        self.lr = learning_rate

    def fit(self, X, y):  # Method Overriding
        print(f"Training {self.model_name} with lr={self.lr}...")

model = LinearRegressionModel(learning_rate=0.05)
model.fit([1, 2], [2, 4])
\end{lstlisting}
\end{examplebox}

\textbf{Data Science Relevance:} Every major Machine Learning framework in Python—including PyTorch (\texttt{nn.Module}), scikit-learn (\texttt{BaseEstimator}), and TensorFlow (\texttt{tf.keras.Model})—is built entirely around OOP class inheritance!

\begin{videocard}{IITM BS Lecture: Inheritance and Method Overriding}{https://www.youtube.com/watch?v=RrFvcDjJptY}
Official lecture demonstrating parent-child class structures, \texttt{super()}, and method overriding.
\end{videocard}

\newpage
\section{Practice Exercises}
\textit{Detailed step-by-step solutions are available in the Appendix.}

\begin{enumerate}
    \item \textbf{Basic Class Creation:} Design a class \texttt{Circle} with an attribute \texttt{radius}. Include methods \texttt{area()} ($A = \pi r^2$) and \texttt{circumference()} ($C = 2 \pi r$). Test with \texttt{r = 5.0}.

    \item \textbf{Class Encapsulation:} Create a \texttt{BankAccount} class with a private attribute \texttt{\_\_balance}. Include methods \texttt{deposit(amount)}, \texttt{withdraw(amount)} (checking for sufficient funds), and \texttt{get\_balance()}.

    \item \textbf{Class vs Instance Variables:} Explain the difference between a class variable and an instance variable in Python. Provide a 5-line code snippet demonstrating a class variable \texttt{species = "Canis lupus"} shared across \texttt{Dog} instances.

    \item \textbf{Inheritance Hierarchy:} Create a parent class \texttt{Vehicle} with attribute \texttt{brand} and method \texttt{drive()}. Create a child class \texttt{ElectricCar} that inherits from \texttt{Vehicle}, adds attribute \texttt{battery\_capacity}, and overrides \texttt{drive()} to print \texttt{"Driving silently on electricity"}.

    \item \textbf{Data Science Application (Custom Transformer Class):} In scikit-learn, data preprocessing transformers implement \texttt{.fit()} and \texttt{.transform()} methods. Create a class \texttt{StandardScaler} with attributes \texttt{mean} and \texttt{std}. Implement:
    \begin{enumerate}
        \item \texttt{fit(data)}: Calculates and stores the mean and standard deviation of a list of floats.
        \item \texttt{transform(data)}: Returns a new list of $z$-score normalized values $z = \frac{x - \text{mean}}{\text{std}}$.
    \end{enumerate}
\end{enumerate}